% 简化的相机成像管线（ISP），各厂商的步骤顺序与实现不同；宽约110mm。
\begin{tikzpicture}[font=\figurefont\small,
  stage/.style={draw=FigureBlue,fill=FigureBlueFill,rounded corners=1.8pt,
    minimum width=20mm,minimum height=8mm,align=center,inner sep=1.5mm},
  nonlin/.style={stage,draw=FigureYellow,fill=FigureYellowFill},
  outstage/.style={stage,draw=FigureRed,fill=FigureRedFill},
  flow/.style={-{Stealth[length=1.8mm]},draw=BookMuted,line width=0.8pt},
  note/.style={text=BookMuted,font=\figurefont\footnotesize}]
  \node[stage] (raw) at (0,0) {RAW 马赛克\\12--14 位线性};
  \node[stage] (black) at (2.9,0) {黑电平扣除\\坏点与暗角校正};
  \node[stage] (wb) at (5.8,0) {白平衡\\通道增益};
  \node[stage] (dm) at (8.7,0) {去马赛克\\每处三通道};
  \node[stage] (ccm) at (8.7,-1.6) {颜色校正\\$3\times3$ 矩阵};
  \node[stage] (nr) at (5.8,-1.6) {降噪\\锐化};
  \node[nonlin] (tone) at (2.9,-1.6) {色调映射\\sRGB 伽马编码};
  \node[outstage] (jpg) at (0,-1.6) {8 位 sRGB\\JPEG 压缩};
  \draw[flow] (raw) -- (black);
  \draw[flow] (black) -- (wb);
  \draw[flow] (wb) -- (dm);
  \draw[flow] (dm) -- (ccm);
  \draw[flow] (ccm) -- (nr);
  \draw[flow] (nr) -- (tone);
  \draw[flow] (tone) -- (jpg);
  \node[note] at (5.8,0.8) {线性光域：数值与进光量成比例};
  \node[note] at (1.45,-2.45) {非线性显示域};
\end{tikzpicture}
